\section{Introduction}
\label{sec:intro}

Online A/B tests on social networks and marketplaces violate the no-interference assumption: a unit's outcome depends on the treatments of its neighbours. The quantity a platform usually wants is the \emph{global average treatment effect} (GATE), $\tau=\frac1n\sum_i[Y_i(\mathbf 1)-Y_i(\mathbf 0)]$, the difference between mean outcomes if everybody and if nobody were treated. The usual difference in means (DiM) is biased for it. Inverse-probability (Horvitz--Thompson, H\'ajek) estimators of exposure-based contrasts~\citep{aronow2017estimating,ugander2013graph} are unbiased for the GATE under a correct exposure model, but need units whose whole neighbourhood is treated (or untreated); on dense graphs such units barely exist and variance explodes. \emph{Regression-based} estimators~\citep{chin2019regression,han2023model} sidestep this: fit $y_i$ on the unit's treatment and graph-derived exposure features from the experimental data, then predict outcomes under the all-treated and all-control assignments and average the difference. They use every unit and are cheap, but how accurate they can be, what determines their error, and when they fail is only partly understood.

\paragraph{This paper.} We study these questions with controlled simulations on synthetic graphs (Erd\H{o}s--R\'enyi, stochastic block model, Watts--Strogatz, Barab\'asi--Albert, random regular) and two public Facebook graphs (Stanford3, $n=11{,}586$; ego-Facebook, $n=4{,}039$), with synthetic outcomes so that the true $\tau$ is known exactly. The device is a decomposition (Sec.~\ref{sec:theory}) of the error into a \emph{design-leverage} term, set by how far the all-treated/all-control exposures lie from those the design generates, and a \emph{bias} term set by misspecification, with closed forms for the one-hop fraction-of-treated-neighbours regression. Findings:
\begin{enumerate}[leftmargin=*,itemsep=1pt,topsep=2pt]
\item \textbf{Sampling error is predictable from graph and design alone.} $\mathrm{RMSE}\approx\hat\kappa\sqrt{L/n}$ with leverage $L=(1+d_H)/(p(1-p))$ under unit randomization ($d_H$: harmonic-mean degree): error grows like $\sqrt{\text{degree}}$, shrinks like $n^{-1/2}$, and one fitted constant predicts 162 Bernoulli-design runs with median ratio $1.01$ ($98\%$ within $\pm10\%$). Cluster randomization lowers $L$ by an amount that depends on how well the graph clusters (RMSE $2.9\times$ lower on ego-Facebook, $1.9\times$ on Stanford3); the closed form is a rough guide there.
\item \textbf{Four bias mechanisms with closed-form predictions that match simulation:} extrapolation under nonlinear exposure response (the regression recovers the slope at $p$, not the secant; Prop.~\ref{prop:extrap}); omitted multi-hop spillover ($B\approx-\beta\lambda^2/(1-\lambda)$; Prop.~\ref{prop:mh}); effect heterogeneity (OLS targets a $1/d_i$-weighted average of unit effects; Prop.~\ref{prop:wt}; bias up to $100\%$ of the GATE on Stanford3); exposure mis-definition. None shrinks with $n$; predicted bias and leverage jointly match the RMSE of 360 Bernoulli world-settings (median ratio $1.005$).
\item \textbf{Design matters as much as the estimator.} Clustering cuts leverage and, where the graph clusters well, bias (multi-hop bias $-0.71\to-0.11$ on ego-Facebook, $-0.85\to-0.49$ on Stanford3, little on ER/SBM); the treated fraction and number of clusters trade bias for variance.
\item \textbf{Rules and negative results.} Regression beats DiM only if spillover exceeds $\approx\hat\kappa\sqrt{L/n}$ ($\ge0.8\beta$ for $n{=}10^3$, $d{=}128$); H\'ajek estimators are undefined on dense random graphs ($d\ge32$) under unit randomization; richer regressions reduce bias only where the graph supplies extreme exposures (a cubic fixes Stanford3/ego-Facebook for $15$--$45\%$ extra RMSE but has RMSE $\approx2$ on near-regular graphs); sandwich intervals cover well when the model is right (mean $0.95$) and fail when it is not (mean $0.57$), the failure being explained by the bias.
\end{enumerate}
\section{Related work}
\label{sec:related}
\textbf{Designs and exposure-based estimation.} \citet{aronow2017estimating} give randomization-based inverse-probability estimators for general exposure mappings; graph cluster randomization~\citep{ugander2013graph}, its randomized variant~\citep{ugander2023randomized} and rate-optimal cluster designs~\citep{leung2022rate} reduce the variance of such estimators; \citet{eckles2017design} study interference bias and neighbourhood-aware analysis; \citet{karrer2021network} describe a deployed system; \citet{chen2023optimized} optimize assignment covariances for HT estimators. \textbf{Outcome models.} \citet{cortez2022staggered,cortez2023exploiting} and \citet{eichhorn2024low} obtain total-effect estimators under low-order interaction models (the last with clustered designs); \citet{yu2022estimating} work without the network. \citet{chin2019regression} proposes the regression-adjustment estimator studied here; \citet{han2023model} select graph- and treatment-based covariates with block bootstrap. \citet{savje2021average} and \citet{li2022random} give asymptotics under limited interference and random graphs. \textbf{This work} asks a diagnostic question about the regression estimator family: what determines its error, and which factor dominates when. Our contributions are closed-form leverage and bias predictors validated against simulation on real graph topologies, a head-to-head with difference-in-means and H\'ajek baselines, and an audit of interval coverage. We propose no new estimator; the weighted and stratified fits in Sec.~\ref{sec:results} are diagnostics-motivated fixes whose merits we only evaluate.

\section{Setting, estimators and metrics}
\label{sec:setting}
\textbf{Setting.} A connected undirected graph with adjacency $A$ on $n$ units, degrees $d_i\ge1$, $W=D^{-1}A$. Assignment $z\in\{0,1\}^n$ is drawn either by \emph{Bernoulli}$(p)$ (independent units) or by \emph{cluster randomization}: units are partitioned into clusters (Louvain communities~\citep{blondel2008fast}, resolution tuned to a target number $K$ of clusters) and clusters are independently treated with probability $p$ (redrawing the rare assignments in which one arm is empty). Potential outcomes $Y_i(z)$ are fixed given a ``world'' (graph plus unit-level parameters); the only randomness within a world is $z$. Exposure features are $s_i=(Wz)_i$ (fraction of treated neighbours), $c_i=(Az)_i$ (count), $s^{(2)}_i=(W^2z)_i$ (two-step fraction).

\textbf{Regression estimators.} $\hat\tau=\frac1n\sum_i[\hat f(x_i(\mathbf 1))-\hat f(x_i(\mathbf 0))]$ where $x_i(z)$ are the features and $\hat f$ is fitted on the observed data. Linear models: $(1,z,s)$ [\textsc{ols-s}, the default, Chin-type]; $(1,z,c)$ [\textsc{ols-cnt}]; $+zs$ [\textsc{int}]; $+s^2$ [\textsc{quad}]; $+s^2,s^3$ [\textsc{cub}]; $+s^{(2)}$ [\textsc{ols-s2}]; $+s\cdot\ell_i,\ \ell_i$ (standardized log degree) [\textsc{s-deg}]; weighted least squares with weights $\propto1/\mathrm{Var}(s_i)$ [\textsc{wls-s}]; separate \textsc{ols-s} fits in five degree strata, averaged with stratum sizes [\textsc{strat-s}]; gradient boosting (scikit-learn HistGradientBoosting) on $(z,s,s^{(2)},\ell)$ [\textsc{gbm}]. \textbf{Baselines:} DiM; H\'ajek estimator with exposure ``unit and all neighbours treated'' / ``none treated'' with exactly computed design probabilities (undefined when no unit qualifies in an arm).

\textbf{Metrics.} For each setting we report, over $R=300$ assignments in each of $6$ independent worlds, the \emph{bias} $\mathbb E\hat\tau-\tau$, \emph{SD}, and \emph{RMSE}$=\sqrt{\mathbb E(\hat\tau-\tau)^2}$; RMSE is our headline metric because it combines bias and variance, which have different drivers; we show both. Tables and figures give means across worlds with min--max across the 6 worlds as error bars (Monte-Carlo standard errors of each world's RMSE are $\le0.04\times$RMSE with $R=300$ and are smaller than the across-world spread). For inference we report empirical coverage of nominal 95\% intervals $\hat\tau\pm1.96\,\widehat{\mathrm{se}}$.
